\chapter{Variante B: Aufbau der Buchstaben}
\label{anhang_aufbau}
Der theoretische Aufbau aller Buchstaben ist hier aufgelistet. Die Buchstaben werden jeweils in Block- und Kursivschrift aufgebaut. Bei Kursivbuchstaben gibt es je nachdem noch Präfixe oder Postfixe, falls mehrere Buchstaben verbunden werden. Diese sind in dieser Auflistung jedoch nicht enthalten.

\section{Kleinbuchstaben}
\begin{table}[!h]
  \begin{center}
    \begin{tabular}{ m{2cm} | c | c }
	\emph{Buchstabe}&\emph{Blockschrift} &\emph{Kursivschrift}\\ \hline
	a &
	\includegraphics[width=15pt]{./img/b_kreis_klein.pdf}
	\includegraphics[width=15pt]{./img/b_gerade_klein.pdf}
	&
	\includegraphics[width=15pt]{./img/b_kreis_klein.pdf}
	\includegraphics[width=15pt]{./img/b_halb_oben.pdf}
	\\ \hline

            b &
	\includegraphics[width=15pt]{./img/a_gerade.pdf}
	\includegraphics[width=15pt]{./img/b_kreis_klein.pdf}	
	&
	\includegraphics[width=15pt]{./img/b_gerade_vor_klein.pdf}
	\includegraphics[width=15pt]{./img/b_kreis_klein.pdf}
	\includegraphics[width=15pt]{./img/b_gerade_klein.pdf}
	\includegraphics[width=15pt]{./img/b_halb_oben.pdf}
	\\ \hline

	 c &
	\includegraphics[width=15pt]{./img/b_halb_rechts.pdf}	
	&
	\includegraphics[width=15pt]{./img/b_halb_rechts.pdf}
	\\ \hline

	 d &
	\includegraphics[width=15pt]{./img/b_kreis_klein.pdf}	
	&
	\includegraphics[width=15pt]{./img/b_halb_rechts.pdf}
	\\ \hline

    \end{tabular}
  \end{center}
  \caption{Kleinbuchstaben Variante B}
  \label{aufbau_buchsatben}
\end{table}